\begin{problem}[理论考试][超导体与磁通涡旋线]{30}
  超导体的一个重要应用是绕制强磁场磁体，其使用的超导线材属于第二类超导体。如果将这类超导体置于磁感应强度为 $B_{a}$ 的外磁场中，其磁力线将以磁通量子（或称为磁通涡旋线）的形式穿透超导体，从而在超导体中形成正三角形的磁通格子，如图1所示。所谓的磁通量子，如图2所示，其中心是半径为 $\xi$ 的正常态（电阻不为零）区域，而其周围处于超导态（电阻为零），存在超导电流，所携带的磁通量为 $\Phi_{0}=\frac{h}{2e}=2.07\times10^{-15}Wb$（磁通量的最小单位）。

  \begin{subproblem}
    若 $B_{a}=5\times10^{-2}T$，求此时磁通涡旋线之间距离 a；
  \end{subproblem}
  \begin{subproblem}
    随着 $B_{a}$ 的增大，磁通涡旋线密度不断增加，当 $B_{a}$ 达到某一临界值 $B_{c2}$ 时，整块超导体都变为正常态。假设磁通涡旋线芯的半径为 $\xi = 5 \times 10^{-9} ~m$，求所对应的临界磁场 $B_{c2}$；
  \end{subproblem}
  \begin{subproblem}
    对于理想的第二类超导体，当有电流 I 通过超导带材时，在安培力的驱动下，磁通涡旋线将会粘滞流动，在超导带内产生电阻（也称为磁通流阻），从而产生焦耳损耗，不利于超导磁体的运行。磁通涡旋线稳定粘滞流动的速度 v 与单位体积磁通涡旋线所受到的驱动力 $f_{A}$ 和 $B_{a}$ 的关系为 $f_{A} = \eta \frac{B_{a}}{\Phi_{0}} v$，其中 $\eta$ 为比例系数。外加磁场、电流方向，以及超导带材的尺寸如图 3 所示，请指出磁通涡旋线流动的方向，并求出磁通涡旋线流动所产生的电阻率（用 $B_{a}$，$\Phi_{0}$，$\eta$，超导体尺寸 b，c，d 表示）；
  \end{subproblem}
  \begin{subproblem}
    要使超导材料真正实用化，消除这种磁通流阻成了技术的关键，请给出你的解决方案。
  \end{subproblem}
\end{problem}